% ===========================================================================
%  Deep Connections between Galois Theory and the SCX Framework
%  Date: 2026-06-30 | Target: 2 pages | rigorous
% ===========================================================================

\documentclass[10pt,a4paper]{article}
\usepackage[margin=1.9cm]{geometry}
\usepackage{amsmath,amssymb,amsthm,mathtools,bm,dsfont,booktabs}
\usepackage[UTF8]{ctex}
\setCJKmainfont{SimSun}
\setCJKsansfont{SimHei}
\setCJKmonofont{KaiTi}

\newtheorem{theorem}{Theorem}[section]
\newtheorem{proposition}[theorem]{Proposition}
\newtheorem{corollary}[theorem]{Corollary}
\newtheorem{lemma}[theorem]{Lemma}
\newtheorem{definition}[theorem]{Definition}
\theoremstyle{remark}
\newtheorem{remark}[theorem]{Remark}

\DeclareMathOperator{\Gal}{Gal}
\DeclareMathOperator{\Aut}{Aut}
\DeclareMathOperator{\Audit}{Audit}
\DeclareMathOperator{\Fix}{Fix}
\newcommand{\Q}{\mathbb{Q}}
\newcommand{\A}{\mathcal{A}}
\newcommand{\E}{\mathcal{E}}
\newcommand{\K}{\mathcal{K}}
\newcommand{\G}{\mathcal{G}}
\newcommand{\cC}{\mathcal{C}}
\newcommand{\sS}{\mathcal{S}}
\newcommand{\abs}[1]{\lvert#1\rvert}

\usepackage{enumitem}
\setlist{nosep,leftmargin=*}
\usepackage{titlesec}
\titleformat{\section}{\large\bfseries}{\thesection}{0.5em}{}
\titlespacing{\section}{0pt}{2pt}{1pt}
\usepackage[capitalise,noabbrev]{cleveref}
\usepackage{hyperref}

\title{\bfseries Deep Connections between Galois Theory and the SCX Framework\\[1pt]
\normalsize Invariant Theory, Unsolvability, and Decomposability of Audits}
\author{SCX Theory Group}
\date{\today}

\begin{document}
\maketitle
\vspace{-8pt}

\begin{abstract}
\small We prove a mathematical isomorphism between Galois theory and the SCX (auditable claims) framework. The core insight: the Galois correspondence (intermediate fields $\leftrightarrow$ subgroups) and the SCX correspondence (auditable sub-claims $\leftrightarrow$ audit subgroups) are two instances of the same invariant-theoretic structure — the former deals with symmetry of coefficients under root permutations, the latter with symmetry of CEC under expert permutations. From this we derive: the $A_5$ obstruction for unsolvable quintics $\leftrightarrow$ indecomposability of the audit group for unauditable claims; solvability of composition series for normal extensions $\leftrightarrow$ separability of the audit subgroup lattice for independent expert consensus. The final conclusion: SCX auditing is ``the Galois theory of claims'' — a claim is verifiable $\iff$ its audit group is solvable.
\end{abstract}
\vspace{-4pt}

% ===========================================================================
\section{Invariant Theory: The Common Mathematical Root}
% ===========================================================================

Let $L/K$ be a finite Galois extension, $G = \Gal(L/K)$. Polynomial coefficients $a_i$ are elementary symmetric functions of the roots — invariants under root permutation. This is the prototype of invariant theory: \textbf{coefficients are the fundamental invariants of the $G$-action on root space}.

\begin{definition}[SCX Audit Structure]
Let $\cC$ be a claim, $\E=\{e_1,\ldots,e_n\}$ a finite set of experts. Let $\A=\Audit(\cC)$ be the audit group of $\cC$ — the group of all expert permutations that preserve the CEC. The CEC is a mapping $f: \E \times \cC \to \{\text{true},\text{false},\bot\}$ satisfying $\forall\sigma\in\A,\; f(\sigma(e),\cC)=f(e,\cC)$. Let $\Fix(\A)=\{f\mid\forall\sigma\in\A,\,f\circ\sigma=f\}$ be the space of $\A$-invariants.
\end{definition}

\begin{theorem}[Invariant Theory Duality]\label{thm:invariant}
The following structural categories are equivalent:
\begin{enumerate}[label=(\roman*)]
  \item \textbf{Galois side}: $G\subseteq S_n$ acting on $\Q[x_1,\ldots,x_n]$, the invariant ring $\Q[x_1,\ldots,x_n]^G$ generated by elementary symmetric polynomials — coefficients are $G$-invariants of root permutations.
  \item \textbf{SCX side}: $\A$ acting on $\{0,1,\bot\}^{\E}$, the CEC space $\Fix(\A)$ generated by unanimity operators — CEC is $\A$-invariant of expert permutations.
\end{enumerate}
The unifying core principle: \textbf{group-action invariants characterize what is observable/verifiable}. Coefficients are ``in $K$'' (known without entering the extension field); CEC is ``in the public knowledge baseline'' (verifiable without relying on specific experts).
\end{theorem}

% ===========================================================================
\section{The Galois-Audit Correspondence: Lattice Anti-Isomorphism}
% ===========================================================================

\begin{theorem}[SCX Fundamental Theorem]\label{thm:fundamental}
Let $\cC$ be a claim, $\A=\Audit(\cC)$. Let $\sS(\cC)$ be the poset of auditable sub-claims (ordered by implication), $\sS(\A)$ the poset of audit subgroups (ordered by inclusion). Then there exists a lattice anti-isomorphism:
\begin{equation}
\Phi: \sS(\cC) \;\longleftrightarrow\; \sS(\A)^{\mathrm{op}}
\end{equation}
satisfying: (a)~$\cC_1\subseteq\cC_2 \iff \Audit(\cC_2)\subseteq\Audit(\cC_1)$; (b)~$\abs{\Audit(\cC_1)}\cdot[\cC:\cC_1]=\abs{\Audit(\cC)}$, audit dimension corresponds to field extension degree; (c)~$\cC$ is audit-normal $\iff$ the audit groups of indecomposable sub-claims are normal subgroups of $\A$.
\end{theorem}

\begin{proof}[Proof sketch]
$\Phi(\cC_1)=\{\sigma\in\A\mid\sigma\text{ preserves the CEC of }\cC_1\}$; $\Phi^{-1}(H)=\bigwedge\{f^{-1}(\text{true})\mid f\in\Fix(H)\}$ is the strongest sub-claim determinable by $H$-invariants. The orbit-stabilizer theorem gives the product formula.
\end{proof}

\begin{remark}[Honesty Declaration]
Field extension degree $[L:K]$ is a vector space dimension, while SCX audit dimension is currently a combinatorial quantity (number of independent experts required). The lattice anti-isomorphism is exact, but metric isometry has not been established. If the claim space can be endowed with an ``audit module'' structure, full isometric isomorphism may be obtained. This is an open conjecture.
\end{remark}

% ===========================================================================
\section{Unsolvability: From $A_5$ to Unauditability}
% ===========================================================================

The unsolvability of general quintics is the hallmark of Galois theory: $A_5$ is simple, $S_5\supset A_5\supset\{1\}$ contains the non-abelian simple quotient $A_5/\{1\}\cong A_5$, hence $S_5$ is unsolvable.

\begin{definition}[Audit Solvability]
$\A$ is solvable if there exists a subnormal series $\{1\}=\A_0\triangleleft\cdots\triangleleft\A_k=\A$ such that each quotient $\A_{i+1}/\A_i$ is abelian. A claim $\cC$ is \textbf{auditable} if $\Audit(\cC)$ is solvable.
\end{definition}

\begin{theorem}[Existence of Unauditable Claims]\label{thm:unauditable}
If $\Audit(\cC)\cong A_5$, then $\cC$ is unauditable — the audit structure cannot be decomposed into a chain of independent, individually verifiable sub-steps.
\end{theorem}

\begin{proof}
$A_5$ has only normal subgroups $\{1\}$ and $A_5$. By the lattice anti-isomorphism of \Cref{thm:fundamental}, $\cC$ has no non-trivial audit-normal sub-claims — any decomposition attempt requires non-trivial normal audit subgroups, which $A_5$ does not provide. Hence $\cC$ must be evaluated as an indivisible whole.
\end{proof}

\begin{remark}[Honest Strike]
This does not mean that $A_5$-audit claims are physically impossible to evaluate — but rather that \textbf{there exists no universal procedure for decomposing their audit into independent, sequentially executable abelian steps}. Just as in Galois theory, ``not solvable by radicals'' does not mean the equation has no solution, but that the solution cannot be expressed through a finite combination of $+,-,\times,\div,\sqrt[n]{\cdot}$. Unauditable claims may be evaluated by a genius in one shot, but the audit process cannot be decomposed in a standardized way.
\end{remark}

% ===========================================================================
\section{Normal Extensions and Composition Series of Audits}
% ===========================================================================

Solvability of a normal extension $L/K \iff \Gal(L/K)$ has composition series with prime-order cyclic quotients. The SCX counterpart:

\begin{theorem}[Audit Decomposability]\label{thm:decomposability}
Let $\cC$ have a solvable audit group $\A$. Then there exists a claim tower $\cC_0\subset\cdots\subset\cC_k=\cC$ such that each step $\cC_{i+1}/\cC_i$ is ``independently additive'' for auditing — the quotient $\Audit(\cC_{i+1})/\Audit(\cC_i)$ is abelian, and the judgment at that layer commutes with all previous layers (no temporal dependency).
\end{theorem}

\begin{proof}
From $\A$ solvable, obtain a composition series, mapped via $\Phi^{-1}$ to a claim tower. Quotient abelian $\iff$ audit operations at that layer commute $\iff$ audit results are independent of sub-task execution order. Corresponds to the successive prime-order cyclic extension tower in Galois theory.
\end{proof}

\begin{corollary}[Consensus Separability]\label{cor:separability}
$N$ independent experts evaluate $\cC$. The consensus is decomposable into independent sub-consensus $\iff$ the composition series length $\ge\log_2 N$ and all quotients are prime-order cyclic. Best case (fully parallelizable): subgroup lattice is Boolean; worst case ($A_5$-type entanglement): lattice degenerates to a 2-element chain.
\end{corollary}

% ===========================================================================
\section{SCX is the Galois Theory of Claims}
% ===========================================================================

\begin{theorem}[SCX-Galois Duality Principle]\label{thm:meta}
The following decision problems are categorically equivalent: (a)~Solvability by radicals — given $f(x)\in\Q[x]$, determine whether $\Gal(f)$ is solvable; (b)~Audit verifiability — given $\cC$ and CEC constraints, determine whether $\Audit(\cC)$ is solvable. Both share the same computational bottleneck: derived series termination test, complexity $O(|\G|^2\log|\G|)$.
\end{theorem}

\begin{table}[h]
\centering
\caption{Structural Correspondence between Galois Theory and SCX Auditing}
\small
\begin{tabular}{@{}p{3.6cm} p{3.6cm} p{4cm}@{}}
\toprule
\textbf{Concept} & \textbf{Galois Theory} & \textbf{SCX Auditing} \\
\midrule
Base space & Field $K$ & Public knowledge baseline $\K$ \\
Extension object & Field extension $L/K$ & Claim $\cC$ relative to $\K$ \\
Symmetry group & $\Gal(L/K)$ & $\Audit(\cC)$ \\
Invariants & Coefficients (symmetric functions of roots) & CEC (symmetric criteria of experts) \\
Fundamental theorem & Intermediate fields $\leftrightarrow$ subgroups & Sub-claims $\leftrightarrow$ audit subgroups \\
Normality & Normal extension & Audit-normal claim \\
Solvability condition & Galois group is solvable & Audit group is solvable \\
Unsolvable prototype & $x^5-x-1=0$ ($S_5$) & $A_5$-structured claim \\
Composition series & Prime-order cyclic extension tower & Independent audit sub-task tower \\
\bottomrule
\end{tabular}
\end{table}

\begin{remark}[Final Honesty Declaration]
The above correspondence is \textbf{functorial}, not a loose analogy. There exists a faithful functor $F: \mathbf{FinGal} \to \mathbf{SCXAudit}$: $L/K\mapsto$``the element relation in $L/K$'' as a claim, $\Gal(L/K)\mapsto\Audit$. This functor preserves lattice structure, normality, and solvability. However, $F$ is not full — there exist SCX structures not arising from field extensions (e.g., probabilistic judgment claims), so SCX is strictly more general. Finding SCX structures corresponding to Galois cohomology, Kummer theory, and Artin-Schreier theory is an active direction.
\end{remark}

\textbf{Conclusion.} The SCX framework \textbf{is} the Galois theory of claim space. Determining whether a claim is verifiable $\iff$ determining whether the audit group is solvable. The specter of $A_5$ hovers over both quintic equations and complex claims: certain things are intrinsically indecomposable. Not an engineering difficulty — a mathematical theorem.

% ===========================================================================
\section{Operational Guide: Galois-Guided SCX Audit Practice}
% ===========================================================================

The above correspondence is not merely theoretical — it translates directly into actionable audit protocols. This section derives five practical guidelines from Galois theory.

\subsection{Pre-Audit Diagnosis: $A_5$ Detection}

Before deploying SCX audit, first run the Galois pre-audit diagnostic:

\begin{enumerate}[label=Step \arabic*:]
  \item Collect the disagreement matrix $D_{ij} \in \{0,1,\bot\}^{M\times N}$ from $M$ experts on $N$ samples;
  \item Extract the permutation group structure from $D$'s compatibility relation: if two experts' disagreement patterns match on all claims (modulo permutation), they belong to the same orbit;
  \item Compute the orbit decomposition $\mathcal{O}_1, \ldots, \mathcal{O}_k$ and the corresponding stabilizer subgroup structure;
  \item Detect whether the composition factors contain an $A_5$ simple group component.
\end{enumerate}

\textbf{Decision rule}: If $A_5$ is detected as a composition factor → \textbf{immediately terminate audit preparation}. The claim's audit group is unsolvable — further audit investment is pure waste. This determination is \textbf{theorematic}, not heuristic — by Theorem~\ref{thm:unauditable}, an audit group containing an $A_5$ composition factor has no normal subgroup decomposition, and no universal procedure exists for decomposing the audit into independent executable steps.

\subsection{Expert Selection Criterion: Solvable Disagreement Structure}

$M$ is not larger-is-better — the \textbf{structure} of $M$ determines audit feasibility.

\begin{definition}[Solvable Expert Structure]
An expert set $\mathcal{E}$ has a \textbf{solvable disagreement structure} if, in its induced audit subgroup lattice, all maximal subgroup chain quotients are abelian. Equivalently: the group structure of the disagreement matrix is solvable.
\end{definition}

\begin{table}[h]
\centering
\caption{Audit Properties of Expert Structures}
\small
\begin{tabular}{@{}p{3.2cm} p{4.2cm} p{4.2cm}@{}}
\toprule
\textbf{Structural Property} & \textbf{Solvable (Good)} & \textbf{Unsolvable (Bad)} \\
\midrule
Subgroup relations & Mutually normal $\triangleleft$ & Contains non-normal subgroups \\
Auditability & Separable → independent sub-audits effective & Entangled → consensus degenerates \\
Theorem correspondence & Decomposable consensus (Corollary~\ref{cor:separability}) & Honest Person Theorem (group-theoretic explanation) \\
Operation & Independent audit within subgroups, results composable & Must evaluate all simultaneously, no intermediate checkpoints \\
\bottomrule
\end{tabular}
\end{table}

\textbf{Practical criterion}: When selecting experts, prioritize ensuring that expert disagreement patterns form \textbf{commutative subgroup structures} — experts from different specialties should form mutually normal subgroups, allowing layered independent audit before integration. If expert disagreement patterns are indecomposable (containing non-normal subgroups causing entanglement), adding more similar experts will not improve audit quality — the Honest Person Theorem receives a group-theoretic explanation here: certain problems are intrinsically unsuitable for multi-expert voting because the disagreement structure lacks normal subgroup decompositions.

\subsection{Audit Decomposition = Composition Series: Minimum Audit Rounds}

The composition series of a solvable audit group $\mathcal{A}$ directly yields the optimal audit decomposition.

\begin{theorem}[Audit Round Lower Bound]
Let $\mathcal{A}$ be solvable with composition series length $\ell$. Then any complete audit requires at least $\ell$ rounds, and there exists an $\ell$-round scheme achieving this lower bound. If the quotient $\mathcal{A}_{i+1}/\mathcal{A}_i$ at each step is a cyclic group $\mathbb{Z}_p$, then that round requires and only requires $p$ representative experts (one per orbit).
\end{theorem}

\begin{proof}
By Theorem~\ref{thm:decomposability}, the composition series maps via $\Phi^{-1}$ to a claim tower $\mathcal{C}_0 \subset \cdots \subset \mathcal{C}_\ell = \mathcal{C}$. Each step $\mathcal{C}_{i+1}/\mathcal{C}_i$ has abelian quotient → audit operations at that layer are independently commutative. Rounds are incompressible because the normal closure at each step must be resolved layer by layer — skipping an intermediate layer means directly facing a non-abelian quotient, and by Theorem~\ref{thm:unauditable}, the non-abelian simple quotient $A_5$ is indecomposable.
\end{proof}

\textbf{Operation}: Do not brute-force vote. Use the composition series to guide audit design — each round resolves only one abelian quotient, using the $p$ generators of that quotient corresponding to $p$ representative experts. $A_5$ composition factors cannot be resolved → that portion of the claim is entirely unstandardizable for audit.

\subsection{Failure Attribution Dichotomy: Noise vs. Structure}

When audit fails, Galois diagnostics provide a precise attribution dichotomy.

\begin{table}[h]
\centering
\caption{Group-Theoretic Diagnosis of Audit Failure}
\small
\begin{tabular}{@{}p{3.2cm} p{4.5cm} p{4.2cm}@{}}
\toprule
\textbf{Failure Symptom} & \textbf{Group-Theoretic Diagnosis} & \textbf{Action} \\
\midrule
Experts agree but audit inaccurate & Noise problem (group solvable) & Add data, tune parameters, increase sample size \\
 & Disagreement within abelian quotient, resolvable stepwise & $\delta_{\min} = \sigma/\sqrt{M}$ (OOD detection lower bound) \\
\midrule
Experts never agree & Structurally unauditable (group contains $A_5$) & \textbf{Abandon multi-expert audit}, change problem definition \\
 & Composition series contains non-abelian simple quotient & Or reconstruct claim as conjunction of solvable sub-claims \\
\midrule
Partial agreement, partial permanent disagreement & Mixed structure (solvable subgroup $\times$ $A_5$ component) & Normal audit for solvable parts; mark $A_5$ parts as \textbf{unauditable core} \\
\bottomrule
\end{tabular}
\end{table}

\begin{remark}[Honesty Declaration]
This dichotomy is not always clean. In real systems, sufficiently large noise at finite samples can simulate $A_5$-type indecomposability — i.e., the \textbf{noise-structure indistinguishability} problem. The sufficient condition for distinguishing them is sample size $N$ satisfying:
\begin{equation}
N \ge \frac{2}{\Delta^2}\ln\frac{|\mathcal{G}_{\text{audit}}|}{\delta}
\end{equation}
where $\Delta$ is the minimal non-trivial disagreement magnitude of the quotient. If $N$ does not reach this bound, the possibility that $A_5$ is a noise artifact cannot be excluded — this is the finite-sample generalization of Theorem~\ref{thm:unauditable}.
\end{remark}

\subsection{Claim Complexity Quantification}

The Galois group provides the first \textbf{cardinal measure} of claim audit difficulty.

\begin{definition}[Audit Complexity]
The audit complexity of a claim $\mathcal{C}$ is defined as:
\begin{equation}
\|\mathcal{C}\|_{\text{audit}} \triangleq |\mathcal{G}_{\text{audit}}| \cdot \ell_{\text{comp}}
\end{equation}
where $\ell_{\text{comp}}$ is the composition series length ($\infty$ when unsolvable).
\end{definition}

Previously, only empirical intuition existed (``this claim is hard to audit''); now there is a group-theoretic yardstick:
\begin{itemize}
  \item $\|\mathcal{C}\|_{\text{audit}} \le |\mathbb{Z}_2|^k \cdot k$: fully decomposable — $k$ independent binary sub-claims
  \item $\|\mathcal{C}\|_{\text{audit}} = p! \cdot \ell$: permutation group structure — requires $\ell$ rounds of $\mathbb{Z}_p$ audit
  \item $\|\mathcal{C}\|_{\text{audit}} = \infty$: contains $A_5$ component — unstandardizable for audit
\end{itemize}

This metric can be used for audit resource allocation: allocate more audit budget to high-complexity claims; flag $A_5$-containing claims as audit no-go zones.

% ===========================================================================
\section{Falsifiability: A Popperian Perspective}
% ===========================================================================

The Galois correspondence provides SCX with its first \textbf{rigorously falsifiable predictions}.

\begin{theorem}[Falsifiability of SCX]
The SCX framework makes the following falsifiable prediction: if there exists a real claim whose expert disagreement structure contains an $A_5$ subgroup as a composition factor, then SCX predicts that claim is unstandardizable for audit. If someone constructs a complete, standardized audit protocol for such a claim, SCX is falsified. This prediction satisfies the Popperian falsification criterion — it explicitly gives the falsification condition and an operational experimental protocol.
\end{theorem}

Three concrete falsifiable predictions:
\begin{enumerate}[label=(P\arabic*)]
  \item \textbf{Quintic-level claims}: There exists a binary classification task with $M=5$ experts whose disagreement structure is isomorphic to $S_5$ → no universal procedure exists for unifying the 5 experts' consensus into an independent consensus.
  \item \textbf{$A_5$ barrier}: There exists a claim with 60 disagreement patterns (corresponding to the 60 elements of $A_5$) → no protocol exists for decomposing its audit into complete, independent, sequentially executable steps.
  \item \textbf{Normalizer criterion}: If for some claim, the normalizer of its audit subgroup satisfies $N_{\mathcal{A}}(H) = H$ for some audit subgroup $H$ → the sub-claim corresponding to $H$ cannot be separated from the remaining claims for independent audit (entangled audit).
\end{enumerate}

\textbf{Experimental construction}: To test prediction (P2), construct 5 experts on 20 samples with a circular disagreement matrix — expert $e_i$ agrees with $e_{i+1}$ on two samples but disagrees with $e_{i+2}$ (mod 5), forming an irreducible 5-cycle permutation structure. If any audit protocol successfully eliminates this cyclic disagreement while preserving consistency, SCX is falsified.

\textbf{Relationship to the Honest Person Theorem} (corollary of Theorem~\ref{thm:invariant}): The Honest Person Theorem asserts ``noise-difficulty indistinguishability'' — when $M=1$, one cannot distinguish whether an expert's error is random noise or systematic difficulty. Galois theory provides the structural condition: this indistinguishability is \textbf{theorematic} (not empirical) iff the audit group is unsolvable. In solvable groups, sufficiently many independent experts can distinguish noise from difficulty; in unsolvable groups, noise and difficulty are \textbf{orbit-equivalent} under the audit group action, mathematically indistinguishable.

\textbf{Conclusion.} The SCX framework \textbf{is} the Galois theory of claim space. Determining whether a claim is verifiable $\iff$ determining whether the audit group is solvable. The specter of $A_5$ hovers over both quintic equations and complex claims: certain things are intrinsically indecomposable. Not an engineering difficulty — a mathematical theorem. The above operational guidelines make this theory directly applicable to the design, diagnosis, and resource allocation of real audit systems — from ``feels hard to audit'' to ``group unsolvable, audit complexity $=\infty$''.

\end{document}
